\documentclass[addpoints]{exam}
\pagestyle{headandfoot}

\usepackage[T1,T2A]{fontenc}
\usepackage[utf8]{inputenc}
\usepackage[russian]{babel}

\usepackage{amsmath, amsfonts}
\usepackage{verbatim}
\usepackage{graphicx}
\usepackage[usenames,dvipsnames]{color}
\usepackage{parskip}
\usepackage{tikz} % drawing schemes
\newcommand{\semester}{WS 2023/2024}
\runningheader{\students}{Submission}{\semester}
\runningfooter{}{\thepage}{}
\headrule
\footrule

% ---------- Modify team name, students (matriculation number), exercise number here ----------
\newcommand{\teamname}{dl2023-team1}
\newcommand{\students}{Jane Doe (1234567), John Doe (7654321)}
\newcommand{\assignmentnumber}{1}
% ---------- End Modify ----------

\title{Submission for Deep Learning Exercise \assignmentnumber}
\author{Team: \teamname\\Students: \students}
\date{\today}

\begin{document}
    \maketitle

    % ---------- Add Solution below here ----------

    \section{Модель Камеры Обскуры}
    
    \begin{enumerate}
    	\item Проецирование. Внутренняя калибровка
    	\item Вращение и перенос. Внешняя калибровка
    \end{enumerate}
    

    \section{Внутренняя калибровка}
    
    
    
    Предположим C и оси X, Y, Z совпадают с камерой,
    ось Z совпадает с принципиальной осью.
    Камера Обскура проецирует 3D точку \( X_0 \) в 2D точку \( X_i \)
    на плоскости изображения, на расстоянии \( f \) от центра проекции.
    
    \begin{tikzpicture}
    	% hint: https://tikz.dev/tutorial
    	\draw[->] (0, 0) -- (6.5, 0) node[anchor=west] {Z}; % axis X
    	\draw[->] (0, 0) -- (0, 2.5) node[anchor=east] {Y}; % axis Y
    	\draw (0, 0) node {\textbullet} node[anchor=east] {$C$} -- (3, 1) node {\textbullet} node[anchor=south] {$X_1$}; % 0 - X1
    	\draw (3, 1) -- (3, 0); % perpendicular
    	\draw (3, 1) -- (6, 2) node[anchor=south] {$X_0$} node {\textbullet}; % X1 - X0
    	\draw (6, 2) -- (6, 0); % perpendicular
    	
    \end{tikzpicture}
    
    У физической Камеры Обскуры плоскость изображения находится после
    центра, в то время как мы ставим плоскость перед центром для простоты.
    
    \[
    \begin{aligned}
    	C &= (0, 0, 0)\\
    	X_0 &= (X, Y, Z)\\
    	X_i &= (x_1, y_1, f)
    \end{aligned}    
    \]
    
    \subsection{Уравнение перспективной проекции}
    \[
    \cfrac{y_1}{f} = \cfrac{Y}{Z};\ \cfrac{x_1}{f} = \cfrac{X}{Z}
    \]
    
    То же самое в однородных координатах:
    
    \[
    \begin{pmatrix}
    	x_1\\
    	y_1\\
    	1
    \end{pmatrix} = 
    \begin{pmatrix}
    	f & 0 & 0 & 0\\
    	0 & f & 0 & 0\\
    	0 & 0 & 1 & 0
    \end{pmatrix}
    \begin{pmatrix}
    	X\\
    	Y\\
    	Z\\
    	1
    \end{pmatrix}
    \]
    
    Получаем
    \[
    x_1 = fX/Z;\ y_1=fY/Z
    \]
    
    \subsection{Преобразование из проекции в пиксельную матрицу}
    
	Нужно учесть размер пикселя \( pix \) и перенести принципиальную точку \( (c_x, c_y) \) в начало координат изображения.
    
    \[
    (x, y) \rightarrow (\frac{x}{pix} + c_x, \frac{y}{pix} + c_y)
    \]
    
    То же самое в виде матриц:
    
    \[
    \begin{pmatrix}
    	u\\
    	v\\
    	1
    \end{pmatrix} =
    \begin{pmatrix}
    	1/pix & s & c_x\\
    	0 & 1/pix & c_y\\
    	0 & 0 & 1
    \end{pmatrix}
    \begin{pmatrix}
    	x\\
    	y\\
    	1
    \end{pmatrix}
    \]
    
    \( s = 0 \) --- наклон осей x и y
    
    
    \subsection{Subsection}

    % ---------- End of Document ----------
    
    \subsection{References}

\end{document}
